\clearpage
\VerdierPhysicalPageAnchor{0148}%
\begingroup
\small
\setlength{\parskip}{0pt}%
\setlength{\emergencystretch}{1.5em}%
\setlength{\abovedisplayskip}{3pt plus 1pt minus 1pt}%
\setlength{\belowdisplayskip}{3pt plus 1pt minus 1pt}%

\begingroup\itshape
\VerdierSemanticAnchor{ch2-2-3-5-a-ii}
\noindent\hspace{1.2em}ii) Tout morphisme $B\longrightarrow X$, avec
$B\in\operatorname{Ob}(\mathcal B)$, $X\in\operatorname{Ob}(\mathcal A)$,
s'insère dans un diagramme commutatif :

\VerdierSemanticAnchor{ch2-2-3-5-a-ii-diagram}
\begin{center}
  \begin{tikzcd}[column sep=4.2em,row sep=2.6em]
    B \arrow[r] \arrow[dr] & X &
    {B'\in\operatorname{Ob}(\mathcal A\cap\mathcal B)} \\
    & B' \arrow[u] &
  \end{tikzcd}
  \quad .
\end{center}

\VerdierSemanticAnchor{ch2-2-3-5-a-definition}
Les sous-catégories triangulées pleines $\mathcal A$ de $\mathcal D$ qui
possèdent les deux propriétés équivalentes ci-dessus sont appelées
\emph{sous-catégories $\mathcal B$-localisantes à droite}.

\VerdierSemanticAnchor{ch2-2-3-5-b}
\par\noindent b) Si $\mathcal A$ est $\mathcal B$-localisante à droite, le
foncteur canonique :

\VerdierSemanticAnchor{ch2-2-3-5-b-functor}
\[
  \mathcal A/(\mathcal A\cap\mathcal B)
  \longrightarrow\mathcal D/\mathcal B
\]
est pleinement fidèle et injectif sur les objets.

\VerdierSemanticAnchor{ch2-2-3-5-c}
\par\noindent c) Les conditions suivantes sont équivalentes :

\VerdierSemanticAnchor{ch2-2-3-5-c-i}
\par\noindent\hspace{1.2em}i) $\mathcal A$ est $\mathcal B$-localisante à
droite.

\VerdierSemanticAnchor{ch2-2-3-5-c-ii}
\par\noindent\hspace{1.2em}ii) Soit
$Q\colon\mathcal D/(\mathcal A\cap\mathcal B)\longrightarrow
\mathcal D/\mathcal B$ le foncteur canonique. La sous-catégorie
$\mathcal A/(\mathcal A\cap\mathcal B)\longrightarrow
\mathcal D/(\mathcal A\cap\mathcal B)$ est formée d'objets $Q$-libres à
droite.
\endgroup

\VerdierSemanticAnchor{ch2-2-3-5-proof-a}
L'équivalence de a) résulte immédiatement de TRIII et des définitions des
systèmes $S_{\mathcal D}(\mathcal B)$ et
$S_{\mathcal A}(\mathcal A\cap\mathcal B)$ (2.1.7).

\VerdierSemanticAnchor{ch2-2-3-5-proof-b-setup}
Démontrons b). Soient $X$ et $Y$ deux objets de $\mathcal A$ (donc de
$\mathcal A/(\mathcal A\cap\mathcal B)$).

\VerdierSemanticAnchor{ch2-2-3-5-proof-b-injectivity}
Montrons d'abord que l'homomorphisme :
\[
  \operatorname{Hom}_{\mathcal A/(\mathcal A\cap\mathcal B)}(X,Y)
  \longrightarrow
  \operatorname{Hom}_{\mathcal D/\mathcal B}(X,Y)
\]
est injectif.

\VerdierSemanticAnchor{ch2-2-3-5-proof-b-quotient-functors}
Désignons par
$Q_1\colon\mathcal A\longrightarrow
\mathcal A/(\mathcal A\cap\mathcal B)$ et
$Q_2\colon\mathcal D\longrightarrow\mathcal D/\mathcal B$ les foncteurs
canoniques de passage au quotient.

\VerdierSemanticAnchor{ch2-2-3-5-proof-b-fraction}
Un élément
$m\in\operatorname{Hom}_{\mathcal A/(\mathcal A\cap\mathcal B)}(X,Y)$ peut
toujours s'écrire sous la forme $m=Q_1(s)^{-1}Q_1(f)$,
$s\in S_{\mathcal A}(\mathcal A\cap\mathcal B)$ (2.2.2). Son image dans
$\operatorname{Hom}_{\mathcal D/\mathcal B}(X,Y)$ est alors
$Q_2(s)^{-1}Q_2(f)$.

\VerdierSemanticAnchor{ch2-2-3-5-proof-b-zero-criterion}
D'après (2.2.4, a)), cette image est nulle si et seulement s'il existe un
diagramme commutatif :

\VerdierSemanticAnchor{ch2-2-3-5-proof-b-zero-diagram}
\begin{center}
  \begin{tikzcd}[column sep=3.7em,row sep=2.5em]
    & {} \arrow[r,"u"] & {} \\
    X \arrow[ur,"f"] && Y \arrow[ul,"s"] \arrow[u,"t"'] &
    {t\in S_{\mathcal D}(\mathcal B)}
  \end{tikzcd}
\end{center}
\endgroup
